%
% This is a borrowed LaTeX template file for lecture notes for CS267,
% Applications of Parallel Computing, UCBerkeley EECS Department.
% Now being used for CMU's 10725 Fall 2012 Optimization course
% taught by Geoff Gordon and Ryan Tibshirani.  When preparing 
% LaTeX notes for this class, please use this template.
%
% To familiarize yourself with this template, the body contains
% some examples of its use.  Look them over.  Then you can
% run LaTeX on this file.  After you have LaTeXed this file then
% you can look over the result either by printing it out with
% dvips or using xdvi. "pdflatex template.tex" should also work.
%
\documentclass[twoside]{article}
\usepackage{graphicx}
\usepackage{tikz}
\usepackage{mathpazo}
\usepackage{pgfplots}

\usetikzlibrary{decorations.pathmorphing}


\setlength{\oddsidemargin}{0.25 in}
\setlength{\evensidemargin}{-0.25 in}
\setlength{\topmargin}{-0.6 in}
\setlength{\textwidth}{6.5 in}
\setlength{\textheight}{8.5 in}
\setlength{\headsep}{0.75 in}
\setlength{\parindent}{0 in}
\setlength{\parskip}{0.1 in}
%
% ADD PACKAGES here:
%
\usepackage{amsmath,amsfonts,graphicx}
%
% The following commands set up the lecnum (lecture number)
% counter and make various numbering schemes work relative
% to the lecture number.
%
\newcounter{lecnum}
\renewcommand{\thepage}{\thelecnum-\arabic{page}}
\renewcommand{\thesection}{\thelecnum.\arabic{section}}
\renewcommand{\theequation}{\thelecnum.\arabic{equation}}
\renewcommand{\thefigure}{\thelecnum.\arabic{figure}}
\renewcommand{\thetable}{\thelecnum.\arabic{table}}
%
% The following macro is used to generate the header.
%
\newcommand{\lecture}[5]{
   \pagestyle{myheadings}
   \thispagestyle{plain}
   \newpage
   \setcounter{lecnum}{#1}
   \setcounter{page}{1}
   \noindent
   \begin{center}
   \framebox{
      \vbox{\vspace{2mm}
    \hbox to 6.28in { {\bf Physics 467/667: Thermodynamics/Statistical Physics
	\hfill Spring 2017} }
       \vspace{4mm}
       \hbox to 6.28in { {\Large \hfill Lecture #1: #2  \hfill} }
       \vspace{2mm}
       \hbox to 6.28in { {\it Lecturer: #3 \hfill Scribes: #4} }
      \vspace{2mm}}
   }
   \end{center}
   \markboth{Lecture #1: #2}{Lecture #1: #2}
   \vspace*{4mm}
}
%
% Convention for citations is authors' initials followed by the year.
% For example, to cite a paper by Leighton and Maggs you would type
% \cite{LM89}, and to cite a paper by Strassen you would type \cite{S69}.
% (To avoid bibliography problems, for now we redefine the \cite command.)
% Also commands that create a suitable format for the reference list.
\renewcommand{\cite}[1]{[#1]}
\def\beginrefs{\begin{list}%
        {[\arabic{equation}]}{\usecounter{equation}
         \setlength{\leftmargin}{2.0truecm}\setlength{\labelsep}{0.4truecm}%
         \setlength{\labelwidth}{1.6truecm}}}
\def\endrefs{\end{list}}
\def\bibentry#1{\item[\hbox{[#1]}]}
%Use this command for a figure; it puts a figure in wherever you want it.
%usage: \fig{NUMBER}{SPACE-IN-INCHES}{CAPTION}
\newcommand{\fig}[3]{
			\vspace{#2}
			\begin{center}
			Figure \thelecnum.#1:~#3
			\end{center}
	}
% Use these for theorems, lemmas, proofs, etc.
\newtheorem{theorem}{Theorem}[lecnum]
\newtheorem{lemma}[theorem]{Lemma}
\newtheorem{proposition}[theorem]{Proposition}
\newtheorem{claim}[theorem]{Claim}
\newtheorem{corollary}[theorem]{Corollary}
\newtheorem{definition}[theorem]{Definition}
\newenvironment{proof}{{\bf Proof:}}{\hfill\rule{2mm}{2mm}}
% **** IF YOU WANT TO DEFINE ADDITIONAL MACROS FOR YOURSELF, PUT THEM HERE:
\newcommand\E{\mathbb{E}}
\begin{document}
%FILL IN THE RIGHT INFO.
%\lecture{**LECTURE-NUMBER**}{**DATE**}{**LECTURER**}{**SCRIBE**}
\lecture{7}{The Second Law and Entropy}{Qiang Zhu}{scribe-name1,2,3}
%\footnotetext{These notes are partially based on those of Nigel Mansell.}
% **** YOUR NOTES GO HERE:
% Some general latex examples and examples making use of the
% macros follow.  
%**** IN GENERAL, BE BRIEF. LONG SCRIBE NOTES, NO MATTER HOW WELL WRITTEN,
%**** ARE NEVER READ BY ANYBODY.

\section{Some useful equations} % Don't be this informal in your notes!
\begin{equation} \label{idealgas} PV = nRT = NkT \end{equation}
\begin{equation} \label{Avogadro} N = n \times N_A \end{equation}
\begin{equation} \label{PV-micro} PV = Nm{\overline v_x^2} = NkT\end{equation}
\begin{equation} \label{eqpartition} U_{thermal} = N \cdot f \cdot \frac{1}{2}kT \end{equation}
\begin{equation} \label{1stlaw} \Delta{U} = Q + W \end{equation}
\begin{equation} \label{work3} \Delta W = - P \Delta{V} \end{equation}

\section{Calculate $\Omega$ for a Einstein Solid}
Remember we have done some computer programmings for two-state systems.
A general trend is that it approaches to a large number $N$, $\Omega(N)$ tends to be localized.
If we treat the histogram of $\Omega(N)$ as a continues function (true when $N$ is very large),
it looks like a very smooth curve and $\Omega(N)$ follows some kind of distribution. 

Now let's try to figure out what it is, by taking a Einstein solid as an example. 
\begin{equation} \label{E1}
 \Omega(N, q) = \binom{q+N-1}{q} = \frac{(q+N-1)!}{q!(N-1)!}
\end{equation}

In reality, there are always many more energy units ($q$) than oscillator ($N$), so we assume $q\gg N$.

To make it easier, let's just remove '-1' in eq. \ref{E1},
\begin{equation} \label{E1}
\begin{split}
 \text{ln}\Omega &= \text{ln}(\frac{(q+N)!}{q!N!})\\
         &= \text{ln}(q+N)! - \text{ln}q! - \text{ln}N!\\
   &\approx (q+N)\text{ln}(q+N) - (q+N) - q\text{ln}q - q - N\text{ln}N - N\\
        & = (q+N)\text{ln}(q+N) - q\text{ln}q - N\text{ln}N \\
\end{split}
\end{equation}

Remember we have the assumption of $q\gg N$, namely ($q/N \rightarrow$ 0)
\begin{equation} \label{E1}
       \text{ln}(q+N) = \text{ln}(q(1+\frac{N}{q})\\
                      = \text{ln}q + \frac{N}{q} 
\end{equation}

Therefore, we have 
\begin{equation} \label{E1}
       \text{ln}\Omega = N\text{ln}\frac{N}{q} + N + \frac{N^2}{q}\\
                 \approx N\text{ln}\frac{N}{q} + N 
\end{equation}

If we remove the logarithm sign,
\begin{equation} \label{E1}
       \Omega(N,q) = (\frac{eq}{N})^N
\end{equation}


\section{Calculate $\Omega$ for two Einstein Solids}
Naturnally, we now know the general form of two Einstein Solids model,
\begin{equation} \label{E1}
       \Omega(N_A,q_A,N_B,q_B) = (\frac{eq_A}{N_A})^{N_A} (\frac{eq_B}{N_B})^{N_B}
\end{equation}

For simplicity, let make $N_A = N_B = N $, then
\begin{equation} \label{E2}
       \Omega(N,q_A,q_B) = (\frac{e}{N})^{2N} (q_A q_B)^{N}
\end{equation}

Based on what we plot in the homework, we know that $\Omega$ reaches its maximum value at 
$q_A = q_B = q$/2,
\begin{equation}
       \Omega_\text{max} = (\frac{e}{N})^{2N} (q/2)^{2N}
\end{equation}

Now let's try to calculate the points near $q$/2, say,
\begin{equation}
  q_A = q/2+x, ~~~~~~~~~ q_B=q/2-x.
\end{equation}

Using eq. \ref{E2},
\begin{equation} \label{E3}
       \Omega(N,q,x) = (\frac{e}{N})^{2N} ((\frac{q}{2})^2-x^2)^{N}.
\end{equation}

To simplify it, we get
\begin{equation} \label{E4}
\begin{split}
       \text{ln}((\frac{q}{2})^2-x^2)^{N} 
           &= N[\text{ln}[(\frac{q}{2})^2-x^2]\\
           &= N[\text{ln}(\frac{q}{2})^2 + \text{ln}(1-(\frac{2x}{q})^2))]\\
           &= N[\text{ln}(\frac{q}{2})^2 - (\frac{2x}{q})^2)]\\
           &= N[\text{ln}(\frac{q}{2})^2 - (\frac{2x}{q})^2)]
\end{split}                                          =
\end{equation}
hence we have 
\begin{equation} \label{E4}
  \Omega = \Omega_\text{max} \cdot e^{-N(2x/q)^2}
\end{equation}
This is a typical Gaussian function. A standard version is as follows,
\begin{equation}
  f(x|\mu,\delta^2) = \frac{1}{\sqrt{2\delta^2\pi}} e^\frac{-(x-\mu)^2}{2\delta^2}.
\end{equation}

\begin{tikzpicture}[x=1.5cm, y=1.5cm]

%normal distribution function 'normaltwo'
    \def\normaltwo{\x,{4*1/exp(((\x-3)^2)/2)}}
% input y parameter
    \def\y{3.5}
% this line calculates f(y)
    \def\fy{4*1/exp(((\y-3)^2)/2)}
% Shade orange area underneath curve.
    %\fill [fill=orange!60] (2.6,0) -- plot[domain=0:4.4] (\normaltwo) -- ({\y},0) -- cycle;
% Draw and label normal distribution function
    \draw[color=blue,domain=0:6] plot (\normaltwo) node[right] {};
% Add dashed line dropping down from normal.
    \draw[dashed] ({\y},{\fy}) -- ({\y},0)     node[below] {$\mu+\delta$};
    \draw[dashed] ({6-\y},{\fy}) -- ({6-\y},0) node[below] {$\mu-\delta$};
% Optional: Add axis labels
    \draw (-.2,2.5) node[left] {$f(x)$};
    \draw (3,-.5) node[below] {$x$};
% Optional: Add axes
    \draw[->] (0,0) -- (6.2,0) node[right] {};
    \draw[->] (0,0) -- (0,5) node[above] {};
\end{tikzpicture}
\begin{enumerate}
\item symmetric
%\item unimodal
%\item its second derivetive is equal to 
\item Gaussian width
\end{enumerate}



The multiplicity falls off to 1/e of its maximum when 
\begin{equation}
  N(\frac{2x}{q})^2 = 1 ~~~~~~or~~~~~~ x = \frac{q}{2\sqrt{N}}
\end{equation}

Let's plug in some number, say $N$=\textrm {$10^{20}$}, 
This results tell us,
when two Einstein solids are in thermodynamical equilibrium, any random fluctuation
will be unmensurable. The most-likely macrostates are very locallized.\\
{\bf Exercises}
\begin{enumerate}
\item Problem 2.20. 
\item Problem 2.23
\end{enumerate}

\section{Ideal Gas}
Suppose we have a single gas atom (Ar), with kinetic energy $U$ in a container of volume $V$, what is its corresponding $\Omega$?
Obviously, the possible microstate is propotional to V. In principle, the atom can stay at any place of $V$.
Also, each microstate can be represented as a vector, since it has velocity (more precisely Momentum). 
Therefore
\begin{equation}
\Omega \approx V\cdot V_p
\end{equation}



It appears that both $V$ and $V_p$ somehow relate to very large numbers, would their product go to infinity? 
Foutunately, we have the famous {\bf Heisenberg uncertainty principle}:
\begin{equation} \Delta{x}\Delta{p_x} = h. \end{equation}

For a one-dimensional chain, we define $L$ as the length in real space, $L_p$ as the length in momentum space,
\begin{equation} \Omega_{1D} = \frac{L}{\Delta{x}} \frac{L_p}{\Delta{p_x}} = \frac{LL_p}{h}. \end{equation}
Therefore, its 3D version is,
\begin{equation} \Omega_1 = \frac{VV_p}{h^3}. \end{equation}
Acccordingly, the multiplicity function for an ideal gas of two molecules should be
\begin{equation} 
\Omega_2 = \frac{1}{2}\frac{V^2}{h^6} \times \text{area of $P$ hypersphere}
\end{equation}
if the two molecules are indistinguishable.
The general form for $N$ should be 
\begin{equation} \Omega_N = \frac{1}{N!}\frac{V^N}{h^{3N}} \times \text{area of $P$ hypersphere}. \end{equation}

For $N$=1, how to calculate the area?
Since $U$ depends on the momentum by
\begin{equation}
U = \frac{1}{2}m(v_x^2+v_y^2+v_z^2)=\frac{1}{2m}(p_x^2+p_y^2+p_z^2)
\end{equation}
\begin{equation}
p_x^2+p_y^2+p_z^2 = \sqrt{2mU}
\end{equation}

So the momentum space is the surface of a sphere with radius $\sqrt{2mU}$, namely,
\begin{equation}
\begin{split}
\text{area} = & 2    ~~~~~~~~~~~~~~~~~~~~~~~~\text {($d$=1)}\\
            = & 2\pi r    ~~~~~~~~~~~~~~~~~~~\text {($d$=2)}\\
            = & 4\pi r^2    ~~~~~~~~~~~~~~~~~\text {($d$=3)}\\
            = & ....\\
            = & ....\\
            = & \frac{2\pi^{d/2}}{\Gamma(d/2)} r^{d-1} ~~~~~\text {($d$ in general)}\\
\end{split}
\end{equation}

Therefore, the general $\Omega$ is
\begin{equation} 
\Omega_N = \frac{1}{N!} \frac{V^N}{h^{3N}} \frac{2\pi^{3N/2}}{(3N/2-1)!} \sqrt{2mU}^{3N-1}. 
\end{equation}


\begin{equation}
\Omega(U,V,N) = f(N)V^NU^{3N/2}
\end{equation}
where $f(N)$ is a complicated function of $N$.\\

For two interaction gases,
\begin{equation}
\Omega(U,V,N) = [f(N)]^2(V_AV_B)^N(U_AU_B)^{3N/2}
\end{equation}




\end{document}

